\begin{abstract}
Low-light images suffer from illumination loss, noise, blur, and color
distortion that reduce both visual quality and downstream recognition
performance. This project investigates whether detection-aware supervision
can improve low-light restoration while preserving restoration fidelity,
using DarkIR as the restoration baseline. We first reproduced DarkIR-m on
LOLBlur and implemented a V1 feature-consistency approach based on normalized
Faster R-CNN feature distances. Matched 50,000-step control and task-aware
experiments showed no consistent detection improvement, motivating a direct
detection-loss formulation.

The final Balanced V2 method trains a DarkIR-m network using a combined L1,
LPIPS, and Sobel edge restoration loss. Detection-guided variants additionally
minimize the four Faster R-CNN detection losses---classification, box
regression, objectness, and RPN box regression---from a frozen
ExDark-trained Faster R-CNN ResNet-50 FPN teacher. Detection weights of
0.05, 0.10, and 0.20 are compared against a matched control, with 0.10
selected as the primary model.

On the 1,020-image scene-disjoint LOLBlur validation split, the control
achieves the best restoration quality. Relative to control, the 0.10 model
decreases PSNR from 33.5311 to 33.4139 dB and SSIM from 0.92083 to 0.91929,
while increasing LPIPS from 0.06793 to 0.06976. Under the ExDark-trained
teacher on 1,800 official-validation images, the 0.10 model improves
mAP@50 from 0.6985 to 0.7189 and mAP@50:95 from 0.4058 to 0.4204 relative
to control. Similar gains occur under an independent frozen
COCO-pretrained detector. Original low-light inputs remain stronger in
absolute detection scores in both protocols. These results demonstrate a
measured restoration--detection trade-off without establishing statistical
significance or universal improvement across raw low-light images.
\end{abstract}